\begin{table}[t]
\centering
\small
\caption{Backtest rejections at 5\% across 8 assets, alpha = 2.5\%}
\label{tab:backtests}
\begin{tabular}{lrrrrr}
\toprule
method & hit\_rate & kupiec & cc & dq & mf \\
\midrule
egarch\_t & 0.018 & 2 & 1 & 0 & 0 \\
ewma & 0.025 & 0 & 0 & 3 & 0 \\
garch\_t & 0.019 & 1 & 1 & 2 & 1 \\
gate & 0.027 & 2 & 1 & 3 & 0 \\
gate\_noexpiry & 0.023 & 1 & 1 & 2 & 0 \\
gate\_noscale & 0.020 & 1 & 1 & 2 & 0 \\
gate\_novix & 0.026 & 1 & 1 & 2 & 0 \\
gate\_qlike & 0.024 & 0 & 0 & 1 & 1 \\
gate\_variance & 0.023 & 0 & 0 & 2 & 2 \\
gjr\_t & 0.018 & 2 & 2 & 0 & 0 \\
har & 0.024 & 0 & 0 & 2 & 0 \\
mean & 0.020 & 1 & 1 & 0 & 0 \\
median & 0.019 & 1 & 1 & 0 & 1 \\
min\_score & 0.022 & 1 & 0 & 1 & 0 \\
previous\_best & 0.022 & 1 & 0 & 2 & 1 \\
relative\_score & 0.021 & 1 & 0 & 0 & 2 \\
vix & 0.022 & 4 & 2 & 3 & 2 \\
\bottomrule
\end{tabular}
\par\footnotesize Counts of assets where each test rejects; hit\_rate is the mean of per-asset hit rates. McNeil-Frey uses residuals scaled by |ES|.
\end{table}
